% Options for packages loaded elsewhere
\PassOptionsToPackage{unicode}{hyperref}
\PassOptionsToPackage{hyphens}{url}
\PassOptionsToPackage{dvipsnames,svgnames,x11names}{xcolor}
\documentclass[
  10pt,
  fontset=fandol]{ctexart}
\usepackage{xcolor}
\usepackage[margin=16mm]{geometry}
\usepackage{amsmath,amssymb}
\setcounter{secnumdepth}{-\maxdimen} % remove section numbering
\usepackage{iftex}
\ifPDFTeX
  \usepackage[T1]{fontenc}
  \usepackage[utf8]{inputenc}
  \usepackage{textcomp} % provide euro and other symbols
\else % if luatex or xetex
  \usepackage{unicode-math} % this also loads fontspec
  \defaultfontfeatures{Scale=MatchLowercase}
  \defaultfontfeatures[\rmfamily]{Ligatures=TeX,Scale=1}
\fi
\usepackage{lmodern}
\ifPDFTeX\else
  % xetex/luatex font selection
\fi
% Use upquote if available, for straight quotes in verbatim environments
\IfFileExists{upquote.sty}{\usepackage{upquote}}{}
\IfFileExists{microtype.sty}{% use microtype if available
  \usepackage[]{microtype}
  \UseMicrotypeSet[protrusion]{basicmath} % disable protrusion for tt fonts
}{}
\makeatletter
\@ifundefined{KOMAClassName}{% if non-KOMA class
  \IfFileExists{parskip.sty}{%
    \usepackage{parskip}
  }{% else
    \setlength{\parindent}{0pt}
    \setlength{\parskip}{6pt plus 2pt minus 1pt}}
}{% if KOMA class
  \KOMAoptions{parskip=half}}
\makeatother
\setlength{\emergencystretch}{3em} % prevent overfull lines
\providecommand{\tightlist}{%
  \setlength{\itemsep}{0pt}\setlength{\parskip}{0pt}}
\usepackage{amsmath,amssymb,mathtools,needspace}
\xeCJKDeclareCharClass{CJK}{"2032,"0400->"04FF,"2160->"216B,"2460->"2473,"25A0,"25C7,"25CB,"25CF}
\IfFontExistsTF{Arial Unicode MS}{\xeCJKsetup{AutoFallBack=true}\setCJKfallbackfamilyfont{\CJKrmdefault}{Arial Unicode MS}}{}
\setcounter{secnumdepth}{0}
\setlength{\emergencystretch}{2em}
\allowdisplaybreaks
\usepackage{bookmark}
\IfFileExists{xurl.sty}{\usepackage{xurl}}{} % add URL line breaks if available
\urlstyle{same}
\hypersetup{
  pdftitle={二阶 Runge-Kutta 方法},
  colorlinks=true,
  linkcolor={Maroon},
  filecolor={Maroon},
  citecolor={Blue},
  urlcolor={Blue},
  pdfcreator={LaTeX via pandoc}}

\title{二阶 Runge-Kutta 方法}
\author{}
\date{}

\begin{document}
\maketitle

\subsubsection{5.3.3 二阶 Runge-Kutta 方法}\label{ux4e8cux9636-runge-kutta-ux65b9ux6cd5}

首先推广改进的 Euler 方法，随意考察区间 \((x_n,x_{n+1})\) 内一点

\[
x_{n+p}=x_n+ph,\qquad 0<p\leqslant1,
\]

希望用 \(x_n\) 和 \(x_{n+p}\) 两个点的斜率值 \(K_1\) 和 \(K_2\) 线性组合得到平均斜率 \(K^*\)，即令

\[
y_{n+1}=y_n+h(\lambda_1K_1+\lambda_2K_2),
\]

其中 \(\lambda_1,\lambda_2\) 为待定系数。同改进的 Euler 格式一样，这里仍取 \(K_1=f(x_n,y_n)\)，问题在于该怎样预测 \(x_{n+p}\) 处的斜率值 \(K_2\)。

仿照改进的 Euler 格式，先用 Euler 格式提供 \(y(x_{n+p})\) 的预测值，即

\[
y_{n+p}=y_n+phK_1,
\]

然后再用预测值 \(y_{n+p}\) 通过计算 \(f\) 产生斜率值

\[
K_2=f(x_{n+p},y_{n+p}).
\]

这样设计出的计算格式具有如下形式：

\[
\begin{cases}
y_{n+1}=y_n+h(\lambda_1K_1+\lambda_2K_2),\\
K_1=f(x_n,y_n),\\
K_2=f(x_{n+p},y_n+phK_1).
\end{cases}\tag{5.3.5}
\]

式 (5.3.5) 含有三个待定系数 \(\lambda_1,\lambda_2\) 和 \(p\)，希望适当选取这些系数的值，使得格式具有二阶精度。

根据式 (5.3.3)，二阶 Taylor 格式为

\[
y_{n+1}=y_n+hy'_n+\frac{h^2}{2}y''_n,
\]

或表示为

\[
y_{n+1}=y_n+hf_n+\frac{h^2}{2}(f_x+f\cdot f_y)_n,
\]

其中 \(f_n\) 和 \((f_x+f\cdot f_y)_n\) 的下标 \(n\) 均表示在 \((x_n,y_n)\) 取值。另一方面，有

\href{../../source-images/pdf-129.jpeg}{原页}

\[
K_1=f_n,\qquad K_2=f(x_{n+p},y_n+phK_1)=f_n+ph(f_x+f\cdot f_y)_n+\cdots,
\]

代入式 (5.3.5)，得

\[
y_{n+1}=y_n+(\lambda_1+\lambda_2)hf_n+\lambda_2ph^2(f_x+f\cdot f_y)_n+\cdots.
\]

由此可见，欲使格式 (5.3.5) 具有二阶精度，只要成立

\[
\begin{cases}
\lambda_1+\lambda_2=1,\\
\lambda_2p=\dfrac{1}{2}.
\end{cases}\tag{5.3.6}
\]

满足条件 (5.3.6) 的一族格式 (5.3.5) 统称为\textbf{二阶 Runge-Kutta 格式}。

除了改进的 Euler 格式外，另一种特殊的二阶 Runge-Kutta 格式是所谓\textbf{变形的 Euler 格式}，其形式如下：

\[
\begin{cases}
y_{n+1}=y_n+hK_2,\\
K_1=f(x_n,y_n),\\
K_2=f\left(x_n+\dfrac{h}{2},y_n+\dfrac{h}{2}K_1\right).
\end{cases}
\]

表面上看，变形的 Euler 格式 \(y_{n+1}=y_n+hK_2\) 仅显含一个斜率值 \(K_2\)，但 \(K_2\) 是通过 \(K_1\) 计算出来的，因此每完成一步仍然需要两次计算函数 \(f\) 的值，工作量和改进的 Euler 格式相同。

总之，二阶 Runge-Kutta 方法用多算一次函数值 \(f\) 的办法，避开了二阶 Taylor 级数法所要求的 \(f\) 导数值的计算。在这种意义上可以说，Runge-Kutta 方法实质上是 Taylor 方法的变形。

\begin{center}\rule{0.5\linewidth}{0.5pt}\end{center}

\subsection{相关原页校注（agent 补充）}\label{ux76f8ux5173ux539fux9875ux6821ux6ce8agent-ux8865ux5145}

以下校注来自本节覆盖的原页，可能也涉及同页相邻小节；原文照录与校正说明分开保留。

\subsubsection{原书 PDF 页 129 的校注}\label{ux539fux4e66-pdf-ux9875-129-ux7684ux6821ux6ce8}

\href{../pages/pdf-129.md}{完整原页转录} · \href{../../source-images/pdf-129.jpeg}{原页图像}

校注记录：ch05a-E003。

\begin{quote}
\textbf{校注（agent 补充，原书下标错误）}：5.3.4 首段``再考察一点 \(x_{n+p}=x_n+qh\)''及后段``为了预测点 \(x_{n+p}\) 处的斜率值 \(K_3\)''，原图两处均印为 \(n+p\)，已保留。按同页对 \(K_3=f(x_{n+q},y_{n+q})\) 的定义及式 (5.3.7)，这两处所指节点均应为 \(x_{n+q}\)。已放大原印字确认，非 OCR 误识。
\end{quote}

\subsection{本节覆盖原页的校注记录}\label{ux672cux8282ux8986ux76d6ux539fux9875ux7684ux6821ux6ce8ux8bb0ux5f55}

下列记录按原页关联，含同页相邻小节的校注；计算时同时读取上文校注和完整原页。

\begin{itemize}
\tightlist
\item
  \texttt{ch05a-E003}：原书 PDF 页 129
\end{itemize}

\end{document}
